\section{Threats to Validity}
\label{sec:threats}

\paragraph{Construct validity.}
The incident taxonomy is normalized from a public archive rather than derived from a single canonical benchmark. This creates unavoidable label noise, especially for records in which the public writeup mixes direct protocol failure with ancillary compromise or vague root-cause descriptions. We mitigate this by treating the incident corpus as a screening layer rather than final exploit evidence. Even so, the residual \texttt{other} category remains large, and this limits any fine-grained statistical claim about exact family prevalence.

\paragraph{Internal validity.}
The protocol-verification results rely on repository reports that already embed analyst judgment about exploitability under a narrow pure on-chain threat model. That judgment is deliberate, but it means the study favors precision over recall. Some incidents that were real in practice may be excluded because they required off-chain compromise, exceptional governance leverage, or market conditions not allowed by the study. Our conclusions should therefore be read as claims about replayable pure on-chain exploitability, not about total ecosystem risk.

\paragraph{External validity.}
The repository covers a non-trivial but still bounded protocol set. Five protocols received deep audits and twenty-four protocols were included in the pure on-chain verification study. These protocols are meaningful because they include major lending, AMM, and staking systems, but they do not represent every DeFi design. Results may not generalize cleanly to long-tail deployments, bespoke bridge architectures, or newly deployed protocols with immature code and low liquidity.

\paragraph{Artifact validity.}
Replay engineering remains an intermediate result rather than a final artifact evaluation. Three replay cases contain concrete logic in Foundry tests, but archive RPC execution was not completed in the present shell configuration. This matters because reproduction difficulty is part of the paper's argument. The artifact set can already support methodological conclusions, but a stronger paper will require archive-validated execution logs for at least a subset of the replay cases.
